\documentclass[10pt,a4paper]{amsart}
\usepackage[margin=19mm]{geometry}
\usepackage{amsmath,amssymb,mathtools}
\usepackage[scr=rsfs,cal=euler]{mathalfa}
\usepackage{xfrac}
\usepackage[unicode,hidelinks,bookmarks=false]{hyperref}
\newcommand{\ie}{i.e.\ }
\newcommand{\cf}{cf.\ }
\newcommand{\resp}{respectively\ }
\newcommand{\Subsection}{\subsection}
\providecommand{\nobreakdash}{\nobreak\hbox{-}}
\setlength{\emergencystretch}{2em}
\multlinegap0pt
\newtheorem{theorem}{Theorem}
\title{Digraph Arborescences and Matrix Determinants}
\author{Sayani Ghosh, Bradley S. Meyer}
\date{Source: arXiv:2309.05827v3 (CC BY 4.0)}
\begin{document}
\maketitle

\noindent\emph{Source and licence.} Digraph Arborescences and Matrix Determinants. Version arXiv:2309.05827v3 (CC BY 4.0), distributed by arXiv under the Creative Commons Attribution 4.0 International licence, http://creativecommons.org/licenses/by/4.0/, as recorded in the arXiv licence metadata for this exact version and shown on the abstract page https://arxiv.org/abs/2309.05827v3. Full licence notice: \texttt{licenses/arxiv-2309.05827-CC-BY-4.0.txt}.\par\smallskip

\noindent\textit{Editorial scope.} Counted unit: the article's theorem eq:detapq (source0.tex lines 675--689). The article's complete original proof of it is delivered with it (source0.tex lines 691--800, one \texttt{proof} environment of 110 lines). No mathematics is written, repaired or re-derived: the statement and the proof are the article's own lines, in the article's own notation, and no retained line is substituted or re-flowed.  Retained uncounted as the source-local context the proof needs: lines 62--79.  Nothing was omitted from any retained span, so the package carries no omission record.  No retained line contains a citation, so the wrapper carries no bibliography and no bibliography entry.  No retained line contains a figure environment, an image-inclusion command or drawing code, so no image is delivered and the package declares no asset dependency.  Editorial formatting only, in addition: an AMS \texttt{amsart} preamble that restates the article's own declarations and macros used by the retained lines, namely the environments \texttt{theorem} on separate counters, together with the standard packages those lines need.  The retained environments print in this document as Theorem 1 in order of appearance, whereas the article prints the same items with the numbering of its own full document.  The complete native source of the article as distributed in the arXiv e-print for this exact version is bundled unchanged as \texttt{sources/146871/source0.tex}.
We begin by considering an $n \times n$ matrix $A = \left[a_{ij}\right]$.
The matrix elements $a_{ij}$ are taken to be
\begin{equation}
a_{ij} = 
\begin{cases}
-v_{ij}, & i\ne j, 1 \leq i, j \leq n\\
 \sum_{k = 1}^n v_{kj}, & i = j, 1 \leq i \leq n
\end{cases}
\label{eq:aij}
\end{equation}
The sum of the elements in column $j$ of the matrix $A$ is thus $v_{jj}$.
Any matrix may be written in the form given by Eq. (\ref{eq:aij}).  The
numbers $v_{ij}$ may themselves be sums.  We thus note that, in general, we
may have
\begin{equation}
v_{ij} = \sum_{\ell=1}^{N_{ij}}  u_{ij}^{(\ell)}
\label{eq:uij}
\end{equation}

\begin{theorem}
Consider the matrix $A$ in Eq. (\ref{eq:aij}).
\begin{equation}
det(A^{P, Q}) = \left(-1\right)^{N_Q} \sum_{B: \{(\tilde{q_k})_1 \to p_j\}}
\epsilon(B) W(B)
\label{eq:detapq}
\end{equation}
where $B$ is an arborescence in the matrix digraph of
$A$ except that $\{(\tilde{q_k})_1 \to p_j\}$
indicates $B$ is rooted at each vertex
$\tilde{q}_k \in \tilde{Q}$ with weight 1 and with a path to some
vertex $p_j \in P$, $W(B)$ is the weight of arborescence $B$ such that $W(B) = \prod_{e \in B} w(e)$, and where $\epsilon(B)$ is a factor $\pm 1$, depending
on the cycles in a subgraph that can be derived from $B$.
\label{theorem:reduced}
\end{theorem}

\begin{proof}
Rearranging the set $Q$ to $\tilde{Q}$ requires $N_Q$
exchanges of ${q_k}'s$.
Since each exchange of $q_k$'s
corresponds to an exchange of columns in $A^{P, Q}$, which, in turn, changes
the sign of the resulting determinant,
\begin{equation}
det(A^{P, Q}) = \left(-1\right)^{N_Q} det(A^{P, \tilde{Q}})
\label{eq:detapqp}
\end{equation}

$A^{P, \tilde{Q}}$ has all zeros in each column $\tilde{q}_k \in \tilde{Q}$
except a 1 in row
$p_k$ for column $\tilde{q}_k$.
If $p_k = \tilde{q}_k$, the matrix digraph for $A^{P, \tilde{Q}}$
has an arc $(0, \tilde{q}_k)$ with weight 1 as the only in arc
to the vertex $\tilde{q}_k$.  Since $p_k = \tilde{q}_k$, $p_k$ and
$\tilde{q}_k$ are the same vertex.  Nevertheless, in this case we consider
there to be a path $\tilde{q}_k \to p_k$.
Let $R \subseteq \tilde{Q}$ be the subset of vertices for which
$\tilde{q}_k = p_k$ for $\tilde{q}_k \in \tilde{Q}$ and $p_k \in P$.

If $p_k \ne \tilde{q}_k$, the vertex $\tilde{q}_k$
will have two in arcs in the $A^{P, \tilde{Q}}$ matrix digraph.
The first will be an arc $(p_k, \tilde{q}_k)$ with weight -1.
Since the diagonal element in $A^{P, \tilde{Q}}$ 
in column $\tilde{q}_k$ and row $\tilde{q}_k$ for this case will be zero,
the second arc is $(0, \tilde{q}_k)$ with weight 1 to cancel out the
weight of the first in arc.
Let $S \subseteq \tilde{Q}$ be the subset of vertices for which
$\tilde{q}_k \ne p_k$ for $\tilde{q}_k \in \tilde{Q}$ and $p_k \in P$.

All arborescences in the matrix digraph of $A^{P, \tilde{Q}}$ must have
either only the arc $(0, \tilde{q}_k)$ if $\tilde{q}_k \in R$ or either
the arc $(0, \tilde{q}_k)$ or $(p_k, \tilde{q}_k)$
if $\tilde{q}_k \in S$.  This means that all arborescences
in the matrix digraph of $A^{P, \tilde{Q}}$ can be derived from an arborescence
rooted at each $\tilde{q}_k \in \tilde{Q}$ by replacing one or more
of the arcs in $(0, \tilde{q}_k)$ for $\tilde{q}_k \in S$ by their
complement arcs $(p_k, \tilde{q}_k)$, as long as the resulting subgraph
does not contain one or more cycles, in which case the subgraph would not
be an arborescence.

Consider an arborescence $B$ in the matrix digraph of $A^{P, \tilde{Q}}$ that
is rooted at each vertex $\tilde{q}_k \in \tilde{Q}$.  Suppose further that
for one of these root vertices $\tilde{q}_j$ there is no path to any
$p_\ell \in P$.  The vertex $p_j$ is thus part of a path from some other
rooted vertex of $B$ to $p_j$.  There is one and only one arborescence
 $B'$ in the
sum over arborescences of the matrix digraph of $A^{P, \tilde{Q}}$ that is
identical to $B$ except that it has the arc $(p_j, \tilde{q}_j)$ with weight
-1 in place of the arc $(0, \tilde{q}_j)$ with weight 1; thus,
$W(B') = -W(B)$, and these arborescences will cancel in the sum over
arborescences.
Since $\vert \tilde{Q} \vert = \vert P \vert$, for an arborescenc $B$ that
is rooted at each vertex $\tilde{q}_k \in \tilde{Q}$, there must
be a path to one and only one vertex $p_j \in P$ from $\tilde{q}_k$.
Recall that we consider
there to be a path $\tilde{q}_k \to p_k$ when $\tilde{q}_k = p_k$.

Now consider an arborescence
$B$ rooted at each vertex $\tilde{q}_k \in \tilde{Q}$
and with a path to one and only one vertex $p_j \in P$.
Consider the subgraphs derived from {\em parent} arborescence
$B$ by replacing one or more
arcs $(0, \tilde{q}_k)$ for $\tilde{q}_k \in S$ by $(p_k, \tilde{q}_k)$.
These subgraphs are arborescences
as long as the replacements do not lead to a cycle or cycles.
Such a cycle $C$ would
be $\tilde{q}_\ell \to ... \to p_j \to \tilde{q}_j \to ... \to p_\ell
\to \tilde{q}_\ell$
and would
result from replacing all $(0, \tilde{q}_k)$ arcs with $(p_k, \tilde{q}_k)$
arcs for each $\tilde{q}_k$ in the cycle.
Let the number of vertices $\tilde{q}_k$ in $C$ be $N_C$.
The sum over arborescence weights derived from parent arborescence
$B$ would be
\begin{equation}
\Sigma(B) = W(B) \prod_{C \in \{C\}_B}
\left(\sum_{r = 0}^{N_C - 1} \binom{N_C}{k}
\left(-1\right)^r 1^{N_C - i}\right)
\label{eq:SigmaB}
\end{equation}
where $C$ is any cycle that can be present among a subset of $\tilde{q}_k$
vertices in $B$ upon replacement of all in arcs to the vertices with
arcs $(p_k, \tilde{q}_k)$ and where $\{C\}_B$ is the set of all such
cycles that can be derived from $B$.  $\Sigma(B)$ may be written
\begin{equation}
\Sigma(B) = W(B) \prod_{C \in \{C\}_B}\left([1 - 1]^{N_C} - \left(-1\right)^{N_C} \right)
= W(B) \prod_{C \in \{C\}_B}\left(-1\right)^{N_C - 1}
\label{eq:SigmaB2}
\end{equation}
Since all arborescences
in the matrix digraph of $A^{P, \tilde{Q}}$ can be derived
from arborescence $B$ rooted at each vertex $\tilde{q}_k \in \tilde{Q}$ with
a path to one and only one $p_j \in P$,
\begin{equation}
\det\left(A^{P, \tilde{Q}}\right) = \sum_{B:\{(q_k)_1 \to p_j\}} \Sigma(B)
\label{eq:detaqp2}
\end{equation}
The arborescences $B$ include only arcs from the digraph of matrix $A$, except
for the root arcs to vertices $\tilde{q}_k$, which have weight 1, so
combination of Eqs. (\ref{eq:detapqp}), (\ref{eq:SigmaB2}),
and (\ref{eq:detaqp2}) with
\begin{equation}
\epsilon(B) = \prod_{C \in \{C\}_B} \left(-1\right)^{\vert C \vert - 1}
\label{eq:epsilonB}
\end{equation}
completes the proof.
\end{proof}

\end{document}
